\section{Partition the experiment without changing its information}
\label{s6:partitioned}
A serial reference is especially useful when calibration, nuisance fitting
and dependent inference interact. In the dynamic funding construction,
workers first accumulate independent audit successes and counts by stratum.
The pilot fitting pass accumulates $X'X$, $X'Y$, response sums and squared
sums on the same fixed training half. Validation losses use the same
candidate penalties and pilot holdout. After choosing the penalty, a second
pilot pass forms the full normal equations. The evaluation pass accumulates
the candidate score numerators and denominators by stratum. No evaluation
outcome chooses the learner or covariance bandwidth.

Let $\hat\psi_t$ be the stacked candidate-response influence, and let
$\bar\psi$ be its global mean. For lag $k$, partition the later index $t$:
\[
 \hat\Gamma_k=T^{-1}\sum_b\sum_{t\in I_b,\,t>k}
   (\hat\psi_t-\bar\psi)(\hat\psi_{t-k}-\bar\psi)'.
\]
Each partition receives a left halo of at most $L$ dates, so cross-boundary
products occur once. Computing a covariance inside each partition and then
averaging would omit those products and could center different experiments.
The resulting Bartlett matrix, inverse-linkage Jacobian, audit covariance
and financial weights are applied in the same order as the serial derivation.
Partition counts are not independent-cluster counts.

Eighteen fresh-process runs compare one, two and four workers at 720,
5,760 and 23,040 evaluation origins, with four times as many pilot origins,
48 branches, seven horizons, two candidate portfolios and two audit strata.
Each worker's numerical libraries are limited to one observed thread.
The measured workflow includes synthetic source generation, revision-aware
materialization, calibration, pilot selection/fitting, response fitting,
covariance and financial result construction. The direct reference comparison,
consumer formatting, semantic checks and receipt serialization are excluded
from the reported workflow time.
All serial/partitioned financial and covariance outputs agree with the
separate reference to below $5\times10^{-10}$ in the recorded absolute
comparisons. Floating-point reduction order can change last bits; it cannot
change the data splits or the number of independent dates.

Mean times over two fresh processes are about .195/.282/.367 seconds at
720 origins and 4.818/5.038/5.024 seconds at 23,040 origins, in
one/two/four-worker order. The largest parent process peaks around 474 MiB;
this whole-process peak also includes the untimed reference comparisons and
semantic checks, rather than isolating the timed workflow's memory.
Child peak-RSS counters are also retained, but are process-accounting maxima,
not a measured simultaneous sum across hosts or workers. Parallel execution
is slower here: process and transfer overhead dominate this small fixed
basis. The evidence supports serial execution as the default for these
dimensions. It does not support a claim of distributed speedup.

The partition construction nevertheless qualifies a harder boundary than
reducing precomputed independent scores: audit uncertainty, pilot model
selection, complete nuisance fits and chronological covariance remain in
one experiment. It is still a single-host, in-memory reference. At scale,
the source materialization and sufficient-statistic passes can be streamed,
but network transfer, object storage, fault recovery and concurrent study
load remain unmeasured. Longer POMDP histories and graph search have
different dominant costs and cannot inherit this timing result. Capacity
planning must combine these measured dimensions with the earlier latent
and discovery workload profiles.
